% ═══════════════════════════════════════════════════════════════════════════════
%  CAPÍTULO — Decorator
% ═══════════════════════════════════════════════════════════════════════════════

\chapter{Decorator}

% ─── Situación Problema ──────────────────────────────────────────────────────
\section{Situación Problema}

Al confirmar una reserva en Airbnb, los huéspedes pueden contratar
servicios adicionales opcionales: seguro de viaje, traslado al aeropuerto,
kit de bienvenida o servicio de limpieza intermedio. Cada combinación
de servicios afecta al precio final y a la descripción del paquete
reservado.

Modelar todas las combinaciones posibles mediante subclases llevaría
a una explosión combinatoria: \texttt{ReservaConSeguro},
\texttt{ReservaConSeguroYTraslado}, \texttt{ReservaConTodosLosServicios},
etc. Con solo cuatro servicios opcionales existen dieciséis combinaciones
posibles. Incorporar un quinto duplicaría ese número y requeriría
modificar la jerarquía de clases cada vez que se añada un servicio nuevo.

% ─── Solución ────────────────────────────────────────────────────────────────
\section{Solución}

El patrón Decorator envuelve el objeto \texttt{Reserva} base con
decoradores concretos: \texttt{DecoradorSeguro},
\texttt{DecoradorTraslado}, \texttt{DecoradorKitBienvenida}. Cada
decorador implementa la misma interfaz que \texttt{Reserva}, delega
las operaciones al objeto que envuelve y añade su propio comportamiento
—incrementar el precio, ampliar la descripción— sin modificar la reserva
base ni los demás decoradores.

Los servicios se pueden combinar libremente en tiempo de ejecución
apilando decoradores en cualquier orden. El código que procesa la
reserva no necesita saber cuántos ni cuáles servicios adicionales se
han contratado; simplemente invoca las operaciones sobre la interfaz
común. Añadir un nuevo servicio es tan sencillo como crear un nuevo
decorador.

% ─── Diagrama de clases ─────────────────────────────────────────────────────
\begin{figure}[H]
    \centering
    % \includegraphics[width=\textwidth]{imagenes/abstract_factory_diagrama.png}
    \caption{Diagrama de clases del patrón Abstract Factory}
\end{figure}


% ─── Roles de las Clases ────────────────────────────────────────────────────
\section{Roles de las Clases}

% Explicar la responsabilidad de cada clase/interfaz

\begin{itemize}
    \item \textbf{AbstractFactory:}
    \item \textbf{ConcreteFactory:}
    \item \textbf{AbstractProduct:}
    \item \textbf{ConcreteProduct:}
    \item \textbf{Client:}
\end{itemize}


% ─── Main ───────────────────────────────────────────────────────────────────
\section{Main}

% Explicar qué realiza el programa principal

\begin{lstlisting}[language=Java, caption=NombreClase]
 
\end{lstlisting}


% ─── Salida del Main ────────────────────────────────────────────────────────
\section{Salida del Main}

\begin{lstlisting}[language=Java, caption=NombreClase]
 
\end{lstlisting}


% ─── Clases ─────────────────────────────────────────────────────────────────
\section{Clases}

% ─── Clase 1 ────────────────────────────────────────────────────────────────
\subsection{NombreClase}

\begin{lstlisting}[language=Java, caption=NombreClase]
 
\end{lstlisting}


% ─── Clase 2 ────────────────────────────────────────────────────────────────
\subsection{NombreClase}

\begin{lstlisting}[language=Java, caption=NombreClase]
 
\end{lstlisting}


% ─── Clase 3 ────────────────────────────────────────────────────────────────
\subsection{NombreClase}

\begin{lstlisting}[language=Java, caption=NombreClase]
 
\end{lstlisting}


% ─── Conclusiones ───────────────────────────────────────────────────────────
\section{Conclusiones}

% Explicar:
% - Ventajas del patrón
% - Cuándo usarlo
% - Beneficios obtenidos
% - Posibles desventajas
% - Relación con principios SOLID